\chapter{Differentiation and Absolute Continuity}\label{lebesgue-differentiation}
\chapterstatus{complete}

\section{Introduction}\label{lebesgue-differentiation:section-introduction}

The fundamental theorem of calculus has two halves: the derivative of $x \mapsto \int_a^xf$ is $f$, and $\int_a^bF^{\prime} = F(b) - F(a)$. This chapter proves both for Lebesgue integrals, with the right
hypotheses, and their generalisations to $\RR^n$ and to measures. The tools are a covering lemma of Vitali and the maximal function of Hardy and Littlewood: Lebesgue's differentiation theorem says that
the averages of an integrable function over small balls around $x$ converge to its value at almost every $x$, and its corollaries are the density theorem and the first half of the fundamental theorem.
The Cantor function shows what can go wrong in the second half. The Radon--Nikodym theorem, proved with von Neumann's Hilbert space argument, says when a measure has a density with respect to another,
and the Lebesgue decomposition splits every measure into an absolutely continuous and a singular part. Finally the absolutely continuous functions are exactly the indefinite integrals, which gives the
fundamental theorem of calculus in its final form. The last section studies functions of bounded variation: Lebesgue's theorem that monotone
functions are differentiable almost everywhere, proved with Vitali's covering theorem for fine covers; Fubini's theorem on series of monotone functions; the decomposition of a monotone
function into absolutely continuous, singular continuous and jump parts; and the theorem of Banach and Zarecki, which characterises absolute continuity by Luzin's property (N). The extension of the Radon--Nikodym theorem to signed and complex measures is in Chapter~\ref{signed-measures}; differentiation of general
Radon measures, and with respect to doubling measures, is in Chapter~\ref{radon-measures}.
References: \cite{folland-real}, \cite{rudin-rca} and \cite{stein-shakarchi-real}.

\noindent\textbf{Prerequisites.} Chapter~\ref{real-analysis} (Real Analysis) and Chapter~\ref{measure-theory} (Measure and Integration).

\section{Differentiation of integrals}\label{lebesgue-differentiation:section-differentiation}

For a continuous function $f$ on $\RR$, the fundamental theorem of calculus says that $\frac{1}{2r}\int_{x-r}^{x+r}f \to f(x)$. We show that the averages of
an integrable function over small balls converge to the function almost everywhere. Let $B(x, r)$ be the open Euclidean ball. The map $y \mapsto x + ry$ multiplies
the volume of boxes by $r^n$, hence Lebesgue measure by $r^n$, so $\lambda(B(x, r)) = r^n\lambda(B(0, 1))$. A measurable function is \emph{locally integrable},
$f \in L^1_{\mathrm{loc}}(\RR^n)$, if $\int_K|f| < \infty$ for every compact $K$.

\begin{lemma}[Vitali covering lemma]\label{lebesgue-differentiation:lemma-vitali-covering}
Every finite family of open balls $B_1, \ldots, B_m$ in $\RR^n$ has a pairwise disjoint subfamily $B_{i_1}, \ldots, B_{i_k}$ with
$\bigcup_iB_i \subseteq \bigcup_j3B_{i_j}$, where $3B(x, r) = B(x, 3r)$. Consequently $\lambda(\bigcup_iB_i) \leq 3^n\sum_j\lambda(B_{i_j})$.
\end{lemma}

\begin{proof}
Choose a ball of largest radius, discard all the balls that meet it, choose a ball of largest radius among the remaining ones, and so on. The chosen balls are
pairwise disjoint. A discarded ball $B = B(c, r)$ meets a chosen ball $B' = B(c', r')$ with $r' \geq r$, namely the one chosen at the step where $B$ was discarded.
If $z \in B \cap B'$ and $y \in B$, then $\|y - c'\| \leq \|y - c\| + \|c - z\| + \|z - c'\| < r + r + r' \leq 3r'$, so $B \subseteq 3B'$. The inequality follows from
$\lambda(3B') = 3^n\lambda(B')$ and subadditivity.
\end{proof}

The proof of Lemma~\ref{lebesgue-differentiation:lemma-vitali-covering} uses only the triangle inequality, so its first statement holds for finite families of balls in every metric space, with
$3B(x, r) = B(x, 3r)$; only the estimate of the measure uses $\RR^n$. For an infinite family the greedy choice may be impossible, because the radii need not attain their supremum. One then chooses,
at each step, a ball whose radius is at least half of the largest available, and the price is the constant $5$ in place of $3$. In a general metric space a ball as a set does not determine its centre
and radius (in $\ZZ$ with the usual distance, $\overline{B}(0, 1/2) = \overline{B}(0, 1/3) = \{0\}$), so balls are given together with a centre and a radius. For a closed ball
$B = \overline{B}(x, r)$ write $\widehat{B} = \overline{B}(x, 5r)$.

\begin{lemma}[$5r$-covering lemma]\label{lebesgue-differentiation:lemma-five-r-covering}
Let $\mathcal{F}$ be a family of closed balls $\overline{B}(x, r)$, $r > 0$, in a metric space, with $\sup\{r\} = R < \infty$. There is a subfamily $\mathcal{G}$ of pairwise disjoint balls such that every
$B \in \mathcal{F}$ meets some $B^{\prime} \in \mathcal{G}$ whose radius is at least half the radius of $B$; in particular $B \subseteq \widehat{B^{\prime}}$.
\end{lemma}

\begin{proof}
Let $\mathcal{F}_j$ be the set of balls in $\mathcal{F}$ with radius in $(2^{-j}R, 2^{1-j}R]$, $j \geq 1$. By Zorn's lemma choose a maximal pairwise disjoint subfamily $\mathcal{G}_1$ of $\mathcal{F}_1$ (the union of a chain of
disjoint subfamilies is disjoint), and inductively a maximal pairwise disjoint subfamily $\mathcal{G}_j$ of the balls in $\mathcal{F}_j$ that are disjoint from all balls of
$\mathcal{G}_1 \cup \cdots \cup \mathcal{G}_{j-1}$. Let $\mathcal{G} = \bigcup_j\mathcal{G}_j$; it is pairwise disjoint. Let $B = \overline{B}(x, r) \in \mathcal{F}_j$. Either $B$ meets a ball of some $\mathcal{G}_i$ with $i < j$,
whose radius exceeds $2^{-i}R \geq 2^{1-j}R \geq r$; or it is disjoint from all of them, and then by maximality of $\mathcal{G}_j$ it meets a ball of $\mathcal{G}_j$, whose radius exceeds $2^{-j}R \geq r/2$. In both
cases $B$ meets $B^{\prime} = \overline{B}(x^{\prime}, r^{\prime}) \in \mathcal{G}$ with $r^{\prime} \geq r/2$. If $z \in B \cap B^{\prime}$ and $y \in B$, then
$d(y, x^{\prime}) \leq d(y, x) + d(x, z) + d(z, x^{\prime}) \leq 2r + r^{\prime} \leq 5r^{\prime}$.
\end{proof}

\begin{counterexample}\label{lebesgue-differentiation:counterexample-five-r-unbounded}
The hypothesis $\sup\{r\} < \infty$ in Lemma~\ref{lebesgue-differentiation:lemma-five-r-covering} cannot be omitted. In $\RR$ let $\mathcal{F} = \{\overline{B}(0, n) : n \geq 1\}$. Any two balls of
$\mathcal{F}$ meet, so a pairwise disjoint subfamily consists of at most one ball $\overline{B}(0, m)$, and then $\overline{B}(0, 5m + 1) \in \mathcal{F}$ is not contained in
$\widehat{\overline{B}(0, m)} = \overline{B}(0, 5m)$. No other constant in place of $5$ helps: the balls $\overline{B}(0, n)$ with $n > Cm$ escape $\overline{B}(0, Cm)$.
\end{counterexample}

\begin{remark}\label{lebesgue-differentiation:remark-five-r-countable}
If the metric space is separable, the subfamily $\mathcal{G}$ of Lemma~\ref{lebesgue-differentiation:lemma-five-r-covering} is countable: each ball $\overline{B}(x, r) \in \mathcal{G}$ contains the
non-empty open ball $B(x, r)$, hence a point of a fixed countable dense set, and disjoint balls contain different points. This applies to $\RR^n$. The proof shows more than the statement: with
$1 < \tau$ in place of $2$ in the dyadic classes (radii in $(\tau^{-j}R, \tau^{1-j}R]$) one gets $r^{\prime} \geq r/\tau$ and hence $B \subseteq \overline{B}(x^{\prime}, (1 + 2\tau)r^{\prime})$, so every
constant larger than $3$ can replace $5$. The constant $3$ itself is reached for finite families (Lemma~\ref{lebesgue-differentiation:lemma-vitali-covering}), where a ball of largest radius exists.
\end{remark}

The $5r$-covering lemma loses a factor $5^n$ in measure, but iterating it recovers everything: an open set can be filled, up to a null set, by disjoint small balls.

\begin{theorem}[Vitali covering theorem]\label{lebesgue-differentiation:theorem-vitali-covering}
Let $U \subseteq \RR^n$ be open with $\lambda(U) < \infty$, and $\delta > 0$. There are countably many pairwise disjoint closed balls $\overline{B}_1, \overline{B}_2, \ldots \subseteq U$ of diameter at most $\delta$
with $\lambda(U \setminus \bigcup_i\overline{B}_i) = 0$.
\end{theorem}

\begin{proof}
\emph{Step 1: a fixed proportion.} Let $V \subseteq \RR^n$ be open with $0 < \lambda(V) < \infty$. We find finitely many disjoint closed balls in $V$ of diameter at most $\delta$ whose union has measure at least
$\theta\lambda(V)$, where $\theta = \frac12 5^{-n}$. Let $\mathcal{F}$ be the family of closed balls contained in $V$ with diameter at most $\delta$; every point of $V$ is the centre of such a ball. By
Lemma~\ref{lebesgue-differentiation:lemma-five-r-covering} and Remark~\ref{lebesgue-differentiation:remark-five-r-countable} there is a countable disjoint subfamily $\mathcal{G}$ with
$V \subseteq \bigcup_{B \in \mathcal{F}}B \subseteq \bigcup_{B^{\prime} \in \mathcal{G}}\widehat{B^{\prime}}$. Hence $\lambda(V) \leq \sum_{\mathcal{G}}5^n\lambda(B^{\prime})$, that is,
$\sum_{\mathcal{G}}\lambda(B^{\prime}) \geq 5^{-n}\lambda(V) > \theta\lambda(V)$, and a finite subfamily already has total measure at least $\theta\lambda(V)$.

\emph{Step 2: iteration.} Put $U_0 = U$. Given the open set $U_k$, if $\lambda(U_k) = 0$ stop; otherwise Step 1 gives finitely many disjoint closed balls in $U_k$ of diameter at most $\delta$, of total measure
at least $\theta\lambda(U_k)$, and $U_{k+1}$ is $U_k$ minus these balls, which is open. Then $\lambda(U_{k+1}) \leq (1 - \theta)\lambda(U_k)$, so $\lambda(U_k) \leq (1 - \theta)^k\lambda(U) \to 0$. All
the balls chosen are disjoint (those chosen at step $k$ lie in $U_k$, which misses the earlier ones), and $U \setminus \bigcup_i\overline{B}_i \subseteq U_k$ for every $k$, so it is a null set.
\end{proof}

The theorem is used in Chapter~\ref{geometric-measure-theory} to identify Hausdorff measure with Lebesgue measure. The proof used only two facts: the $5r$-covering lemma, valid in every metric
space, and $\lambda(\overline{B}(x, 5r)) = 5^n\lambda(\overline{B}(x, r))$. So it works verbatim for balls of any norm on $\RR^n$ (for instance cubes, the balls of the maximum norm), and with
$5^n$ replaced by a power of the doubling constant, for doubling measures on separable metric spaces (Chapter~\ref{radon-measures}).

The theorem fills an open set. Lebesgue's original form of the covering theorem fills an arbitrary set, using only balls from a prescribed family that contains arbitrarily small balls around each of
its points; this is what is needed to differentiate monotone functions in Section~\ref{lebesgue-differentiation:section-bounded-variation}. For an arbitrary $E \subseteq \RR^n$, the outer measure
$\lambda^*(E)$ is the infimum of $\lambda(U)$ over open $U \supseteq E$: enlarge the boxes of a cover of $E$ slightly to open boxes.

\begin{definition}\label{lebesgue-differentiation:definition-fine-cover}
Let $E \subseteq \RR^n$. A family $\mathcal{F}$ of closed balls of positive radius is a \emph{fine cover} (or \emph{Vitali cover}) of $E$ if for every $x \in E$ and every $\delta > 0$ there is a ball $B \in \mathcal{F}$
with $x \in B$ and $\operatorname{diam}B < \delta$.
\end{definition}

The balls need not be centred at the points of $E$. In $\RR$ the closed balls are the closed intervals $[c, d]$, $c < d$; for instance the intervals $[x, x + h]$ and $[x - h, x]$ for $x \in E$ and
$0 < h < 1$ form a fine cover of $E$.

\begin{corollary}[Vitali covering theorem for fine covers]\label{lebesgue-differentiation:corollary-vitali-fine}
Let $E \subseteq \RR^n$ with $\lambda^*(E) < \infty$, and let $\mathcal{F}$ be a fine cover of $E$. There are countably many pairwise disjoint balls $B_1, B_2, \ldots \in \mathcal{F}$ with
$\lambda^*(E \setminus \bigcup_iB_i) = 0$. Consequently, for every $\varepsilon > 0$ there are finitely many pairwise disjoint balls $B_1, \ldots, B_m \in \mathcal{F}$ with
$\lambda^*(E \setminus \bigcup_{i=1}^mB_i) < \varepsilon$, and given an open $U \supseteq E$ the balls can be chosen inside $U$.
\end{corollary}

\begin{proof}
Fix an open $U \supseteq E$ with $\lambda(U) < \infty$, and let $\theta = \frac145^{-n}$. We may assume $\lambda^*(E) > 0$.

\emph{Step 1: a fixed proportion.} Let $E^{\prime} \subseteq E$ with $\lambda^*(E^{\prime}) > 0$, and $U^{\prime} \supseteq E^{\prime}$ open. We find finitely many disjoint balls of $\mathcal{F}$ in $U^{\prime}$ with
$\lambda^*(E^{\prime} \setminus \bigcup B_i) \leq (1 - \theta)\lambda^*(E^{\prime})$. Shrinking $U^{\prime}$, we may assume $\lambda(U^{\prime}) \leq (1 + \theta)\lambda^*(E^{\prime})$. Let $\mathcal{F}^{\prime}$ be the
family of balls of $\mathcal{F}$ of diameter at most $1$, contained in $U^{\prime}$ and meeting $E^{\prime}$; since $U^{\prime}$ is open and $\mathcal{F}$ is fine, every point of $E^{\prime}$ lies in a ball of
$\mathcal{F}^{\prime}$. By Lemma~\ref{lebesgue-differentiation:lemma-five-r-covering} and Remark~\ref{lebesgue-differentiation:remark-five-r-countable} there is a countable disjoint subfamily $\mathcal{G}$ with
$E^{\prime} \subseteq \bigcup_{\mathcal{G}}\widehat{B^{\prime}}$, so $\lambda^*(E^{\prime}) \leq 5^n\sum_{\mathcal{G}}\lambda(B^{\prime})$, and finitely many balls $B_1, \ldots, B_m$ of $\mathcal{G}$ have total measure at
least $\frac125^{-n}\lambda^*(E^{\prime}) = 2\theta\lambda^*(E^{\prime})$. Then
\[
\lambda^*\Bigl(E^{\prime} \setminus \bigcup_iB_i\Bigr) \leq \lambda\Bigl(U^{\prime} \setminus \bigcup_iB_i\Bigr) = \lambda(U^{\prime}) - \sum_i\lambda(B_i) \leq (1 + \theta - 2\theta)\lambda^*(E^{\prime}).
\]

\emph{Step 2: iteration.} Apply Step 1 to $E$ and $U$, then to $E \setminus (B_1 \cup \cdots \cup B_m)$ and the open set $U \setminus (B_1 \cup \cdots \cup B_m)$, and so on, stopping if the remainder becomes null.
The balls are pairwise disjoint and lie in $U$, and after $k$ rounds the remainder has outer measure at most $(1 - \theta)^k\lambda^*(E) \to 0$.

\emph{Step 3: finitely many balls.} The disjoint balls lie in $U$, so $\sum_i\lambda(B_i) \leq \lambda(U) < \infty$; choose $m$ with $\sum_{i > m}\lambda(B_i) < \varepsilon$. Then
$\lambda^*(E \setminus \bigcup_{i \leq m}B_i) \leq \lambda^*(E \setminus \bigcup_iB_i) + \sum_{i > m}\lambda(B_i) < \varepsilon$.
\end{proof}

\begin{definition}\label{lebesgue-differentiation:definition-maximal-function}
For $f \in L^1_{\mathrm{loc}}(\RR^n)$, $x \in \RR^n$ and $r > 0$, the \emph{average} of $f$ over $B(x, r)$ is
$A_rf(x) = \frac{1}{\lambda(B(x, r))}\int_{B(x, r)}f\,d\lambda$, and the \emph{Hardy--Littlewood maximal function} is $Mf(x) = \sup_{r > 0}A_r|f|(x) \in [0, \infty]$.
\end{definition}

\begin{theorem}[Hardy--Littlewood maximal inequality]\label{lebesgue-differentiation:theorem-maximal-inequality}
Let $f \in L^1(\RR^n)$ and $t > 0$. The set $\{Mf > t\}$ is open, and $\lambda(\{Mf > t\}) \leq \frac{3^n}{t}\|f\|_1$.
\end{theorem}

\begin{proof}
If $A_r|f|(x) > t$, choose $s > r$ with $(r/s)^nA_r|f|(x) > t$. For $\|x' - x\| < s - r$ we have $B(x', s) \supseteq B(x, r)$, so
$A_s|f|(x') \geq (r/s)^nA_r|f|(x) > t$; hence $\{Mf > t\}$ is open. Let $K \subseteq \{Mf > t\}$ be compact. Every $x \in K$ is the centre of a ball $B_x$ with
$\int_{B_x}|f| > t\lambda(B_x)$, and finitely many of these balls cover $K$. By Lemma~\ref{lebesgue-differentiation:lemma-vitali-covering} they contain disjoint balls
$B_1, \ldots, B_k$ with
\[
\lambda(K) \leq 3^n\sum_j\lambda(B_j) < \frac{3^n}{t}\sum_j\int_{B_j}|f| \leq \frac{3^n}{t}\|f\|_1.
\]
By inner regularity (Theorem~\ref{measure-theory:theorem-lebesgue-measure}), the same bound holds for $\{Mf > t\}$.
\end{proof}

\begin{theorem}[Lebesgue differentiation theorem]\label{lebesgue-differentiation:theorem-lebesgue-differentiation}
Let $f \in L^1_{\mathrm{loc}}(\RR^n)$. For almost every $x \in \RR^n$,
\[
\lim_{r \to 0}\frac{1}{\lambda(B(x, r))}\int_{B(x, r)}|f(y) - f(x)|\,d\lambda(y) = 0,
\]
and in particular $A_rf(x) \to f(x)$ as $r \to 0$.
\end{theorem}

\begin{proof}
For $x \in B(0, R)$ and $r < 1$ the quantity only involves $f$ on $B(0, R + 1)$, so, replacing $f$ by $f\mathbf{1}_{B(0, R+1)}$ and letting $R$ run through $\NN$, we
may assume $f \in L^1(\RR^n)$. Let $\Omega_f(x) = \limsup_{r \to 0}A_r|f - f(x)|(x)$. Fix $t > 0$ and $\varepsilon > 0$, and choose $g \in C_c(\RR^n)$ with
$\|f - g\|_1 < \varepsilon$ (Proposition~\ref{measure-theory:proposition-density}). For the continuous $g$, $A_r|g - g(x)|(x) \leq \sup_{B(x, r)}|g - g(x)| \to 0$, so
$\Omega_g = 0$. With $h = f - g$ we have $|f(y) - f(x)| \leq |h(y)| + |h(x)| + |g(y) - g(x)|$, hence $\Omega_f(x) \leq Mh(x) + |h(x)|$, and
\[
\{\Omega_f > 2t\} \subseteq \{Mh > t\} \cup \{|h| > t\}.
\]
By Theorem~\ref{lebesgue-differentiation:theorem-maximal-inequality} and Markov's inequality (Corollary~\ref{measure-theory:corollary-integral-properties}(4)), the right
side has measure at most $3^n\varepsilon/t + \varepsilon/t$. As $\varepsilon$ is arbitrary, $\{\Omega_f > 2t\}$ has outer measure $0$, so it is a null set (Lebesgue
measure is complete), and so is the union of these sets for $t = 1/k$, $k \geq 1$. Finally $|A_rf(x) - f(x)| \leq A_r|f - f(x)|(x)$.
\end{proof}

The points where the conclusion holds are the \emph{Lebesgue points} of $f$; changing $f$ on a null set does not change the averages, so almost every point is
a Lebesgue point of every representative.

\begin{corollary}[Lebesgue density theorem]\label{lebesgue-differentiation:corollary-density}
Let $A \subseteq \RR^n$ be Lebesgue measurable. Then $\lambda(A \cap B(x, r))/\lambda(B(x, r)) \to 1$ as $r \to 0$ for almost every $x \in A$, and $\to 0$ for almost
every $x \notin A$.
\end{corollary}

\begin{proof}
Apply Theorem~\ref{lebesgue-differentiation:theorem-lebesgue-differentiation} to the bounded, hence locally integrable, function $\mathbf{1}_A$, whose average over $B(x, r)$ is
the ratio in question.
\end{proof}

\begin{corollary}[Fundamental theorem of calculus for Lebesgue integrals]\label{lebesgue-differentiation:corollary-ftc-lebesgue}
Let $f \in L^1([a, b])$ and $F(x) = \int_{[a, x]}f\,d\lambda$. Then $F$ is continuous, and at almost every $x \in (a, b)$ it is differentiable with $F'(x) = f(x)$.
\end{corollary}

\begin{proof}
Extend $f$ by $0$ to $\RR$. For $y \to x$, $|F(y) - F(x)| \leq \int\mathbf{1}_{[x, y]}|f| \to 0$ by dominated convergence, where $[x, y]$ means the interval between
$x$ and $y$. For $h \neq 0$ let $I_h$ be the closed interval between $x$ and $x + h$, of length $|h|$. Then $F(x + h) - F(x) = \pm\int_{I_h}f$, with the sign of
$h$, so in both cases $\frac{F(x + h) - F(x)}{h} - f(x) = \frac{1}{|h|}\int_{I_h}(f - f(x))$. As $I_h \subseteq \overline{B}(x, |h|)$, an interval of length $2|h|$
whose endpoints are null,
\[
\Bigl|\frac{F(x + h) - F(x)}{h} - f(x)\Bigr| \leq \frac{1}{|h|}\int_{B(x, |h|)}|f - f(x)| = 2A_{|h|}|f - f(x)|(x),
\]
which tends to $0$ at every Lebesgue point of $f$.
\end{proof}

\begin{counterexample}\label{lebesgue-differentiation:counterexample-cantor-function}
\emph{The Cantor function.} A continuous nondecreasing function can have derivative $0$ almost everywhere without being constant, so the formula
$F(b) - F(a) = \int_a^bF'$ fails for such functions. Let $C \subseteq [0, 1]$ be the Cantor set, the set of numbers $\sum_{k \geq 1}a_k3^{-k}$ with all
$a_k \in \{0, 2\}$; these expansions are unique, and $C$ is compact of measure zero, being covered, for each $k$, by $2^k$ intervals of length $3^{-k}$
(Exercise~\ref{measure-theory:exercise-cantor}). Define $c \colon C \to [0, 1]$ by $c(\sum_ka_k3^{-k}) = \sum_k(a_k/2)2^{-k}$.
\begin{enumerate}
\item $c$ is nondecreasing and surjective. If $x < y$ in $C$, their digits first differ at some $k$ with $a_k = 0$ and $b_k = 2$ (a difference $2 \cdot 3^{-k}$ in
position $k$ cannot be compensated by the later digits, which contribute at most $\sum_{j > k}2 \cdot 3^{-j} = 3^{-k}$); then $c(y) - c(x) \geq 2^{-k} - \sum_{j > k}2^{-j} = 0$.
Every $t \in [0, 1]$ has a binary expansion $\sum_kb_k2^{-k}$, and $t = c(\sum_k2b_k3^{-k})$.
\item Let $F(x) = \sup\{c(y) : y \in C, y \leq x\}$ for $x \in [0, 1]$. Then $F$ is nondecreasing, $F = c$ on $C$, and $F([0, 1]) = [0, 1]$. A nondecreasing function
whose image is an interval has no jumps, so $F$ is continuous.
\item The complement $[0, 1] \setminus C$ is a countable union of open intervals $(\alpha, \beta)$ with $\alpha \in C$, removed in the construction; on each of them
$F = c(\alpha)$ is constant, as $C$ has no point in $(\alpha, \beta)$. So $F' = 0$ on $[0, 1] \setminus C$, that is, almost everywhere, while $F(1) - F(0) = 1$.
\end{enumerate}
So $F$ is not the integral of its derivative. The functions that are integrals of their derivatives are exactly the \emph{absolutely continuous} ones
\cite[Chapter 7]{rudin-rca}. Note also that the continuous map $F$ sends the null set $C$ onto $[0, 1]$.
\end{counterexample}

\begin{counterexample}\label{lebesgue-differentiation:counterexample-non-borel}
\emph{A Lebesgue measurable set that is not a Borel set.} First, every $A \subseteq \RR$ with $\lambda^*(A) > 0$ contains a non-measurable subset. Let $V$ be as in
Counterexample~\ref{measure-theory:counterexample-vitali}. If $S$ is a measurable subset of a translate $V + q$, $q \in \QQ$, then the sets $S - q + q'$,
$q' \in \QQ \cap [-1, 1]$, are pairwise disjoint (they lie in the disjoint translates $V + q'$), contained in $[-1, 2]$, and all of measure $\lambda(S)$; so
$\lambda(S) = 0$. The sets $A \cap (V + q)$, $q \in \QQ$, cover $A$, since every real number differs from an element of $V$ by a rational; if they were all
measurable, $A$ would be a countable union of null sets.

Now let $F$ be the Cantor function and $h(x) = x + F(x)$ on $[0, 1]$. It is continuous and strictly increasing from $[0, 1]$ onto $[0, 2]$, hence a
homeomorphism. It maps each interval $(\alpha, \beta)$ of $[0, 1] \setminus C$ onto an interval of the same length, since $F$ is constant there; these images are
disjoint, so $\lambda(h([0, 1] \setminus C)) = 1$ and $\lambda(h(C)) = 1$. Choose a non-measurable $N \subseteq h(C)$. Then $h^{-1}(N) \subseteq C$ is a subset of a
null set, hence Lebesgue measurable. It is not a Borel set: otherwise $N$, its preimage under the continuous map $h^{-1}$, would be Borel, hence measurable.
In particular, a composite of a Lebesgue measurable function with a continuous one need not be Lebesgue measurable: $\mathbf{1}_{h^{-1}(N)} \circ h^{-1} = \mathbf{1}_N$.
\end{counterexample}

As an application of the layer cake formula (Proposition~\ref{measure-theory:proposition-layer-cake}) we show that the maximal function of an $L^p$ function is in $L^p$
for $p > 1$, although it is never in $L^1$ (Exercise~\ref{lebesgue-differentiation:exercise-maximal-not-integrable}).

\begin{corollary}[Hardy--Littlewood maximal theorem]\label{lebesgue-differentiation:corollary-maximal-lp}
For $1 < p \leq \infty$ and $f \in L^p(\RR^n)$, $\|Mf\|_p \leq C_{n,p}\|f\|_p$, with $C_{n,\infty} = 1$ and $C_{n,p}^p = 2^p3^np/(p - 1)$.
\end{corollary}

\begin{proof}
For $p = \infty$, every average of $|f|$ is at most $\|f\|_\infty$. Let $p < \infty$ and $t > 0$, and split $f = f_1 + f_2$ with $f_1 = f\mathbf{1}_{\{|f| > t/2\}}$. Then $f_1 \in L^1$,
since $|f_1| \leq (2/t)^{p-1}|f|^p$, and $Mf \leq Mf_1 + Mf_2 \leq Mf_1 + t/2$, so $\{Mf > t\} \subseteq \{Mf_1 > t/2\}$, and by
Theorem~\ref{lebesgue-differentiation:theorem-maximal-inequality}, $\lambda(\{Mf > t\}) \leq \frac{2 \cdot 3^n}{t}\int_{\{|f| > t/2\}}|f|$. By
Proposition~\ref{measure-theory:proposition-layer-cake} and Tonelli's theorem,
\[\begin{aligned}
\int(Mf)^p &= p\int_0^\infty t^{p-1}\lambda(\{Mf > t\})\,dt \leq 2 \cdot 3^np\int|f(x)|\int_0^{2|f(x)|}t^{p-2}\,dt\,d\lambda(x) \\
&= \frac{2 \cdot 3^np}{p - 1}\int|f|(2|f|)^{p-1},
\end{aligned}\]
which is $\frac{2^p3^np}{p - 1}\|f\|_p^p$.
\end{proof}

\section{The Radon--Nikodym and Riesz representation theorems}\label{lebesgue-differentiation:section-radon-nikodym}

\begin{lemma}\label{lebesgue-differentiation:lemma-hilbert-riesz}
Let $H$ be a complex inner product space that is complete (a \emph{Hilbert space}), for instance $L^2(\mu)$. Every continuous linear functional $\Lambda \colon H \to \CC$ is of the form $\Lambda(f) = \langle f, g \rangle$
for a unique $g \in H$.
\end{lemma}

\begin{proof}
If $\Lambda = 0$ take $g = 0$. Otherwise $N = \ker\Lambda$ is a closed proper subspace. Choose $u \notin N$ and let $d = \inf_{v \in N}\|u - v\| > 0$. For $v_k \in N$ with $\|u - v_k\| \to d$, the parallelogram law
$\|v_k - v_l\|^2 = 2\|u - v_k\|^2 + 2\|u - v_l\|^2 - 4\|u - \frac{v_k + v_l}{2}\|^2$ shows that $(v_k)$ is Cauchy; its limit $v$ satisfies $\|u - v\| = d$. For $w \in N$ and $t \in \CC$,
$\|u - v - tw\|^2 \geq d^2$ forces $\langle u - v, w \rangle = 0$. So $z = u - v \neq 0$ is orthogonal to $N$. For $f \in H$, $\Lambda(f)z - \Lambda(z)f \in N$, so $\Lambda(f)\|z\|^2 = \Lambda(z)\langle f, z \rangle$, and
$g = \overline{\Lambda(z)}z/\|z\|^2$ works. Uniqueness: if $\langle f, g - g' \rangle = 0$ for all $f$, take $f = g - g'$.
\end{proof}

\begin{theorem}[Radon--Nikodym]\label{lebesgue-differentiation:theorem-radon-nikodym}
Let $\mu$ and $\nu$ be $\sigma$-finite measures on $(X, \mathcal{A})$ with $\nu \ll \mu$, meaning $\mu(E) = 0 \Rightarrow \nu(E) = 0$. There is a measurable $h \geq 0$, unique a.e., with
$\nu(E) = \int_Eh\,d\mu$ for all $E$.
\end{theorem}

\begin{proof}
Assume $\mu$ and $\nu$ finite; the $\sigma$-finite case follows by decomposing $X$. Let $\rho = \mu + \nu$. On $L^2(\rho)$, $f \mapsto \int f\,d\nu$ is continuous, since
$|\int f\,d\nu| \leq \int|f|\,d\rho \leq \rho(X)^{1/2}\|f\|_2$. By Lemma~\ref{lebesgue-differentiation:lemma-hilbert-riesz} there is $g$ with $\int f\,d\nu = \int fg\,d\rho$ for all $f \in L^2(\rho)$, and $g$ may be taken real. Taking
$f = \mathbf{1}_E$ gives $\nu(E) = \int_Eg\,d\rho$, so $g \geq 0$ $\rho$-a.e.; and $\int_E(1 - g)\,d\nu = \int_Eg\,d\mu$ (for bounded $f$, $\int f(1 - g)\,d\nu = \int fg\,d\mu$). For $E = \{g \geq 1\}$ the left side
is $\leq 0$ and the right side is $\geq \mu(E)$, so $\mu(E) = 0$ and $\nu(E) = 0$. Applying the identity to $f = \mathbf{1}_E(1 + g + \cdots + g^k)$ gives $\int_E(1 - g^{k+1})\,d\nu = \int_E(g + \cdots + g^{k+1})\,d\mu$. As
$k \to \infty$, monotone convergence gives $\nu(E) = \int_Eh\,d\mu$ with $h = g/(1 - g)$. If $h'$ also works, $\int_E(h - h')\,d\mu = 0$ for $E = \{h > h'\}$ and $\{h < h'\}$ (intersected with sets of finite measure), so
$h = h'$ a.e.
\end{proof}

Two measures $\mu$, $\nu$ on $(X, \mathcal{A})$ are \emph{mutually singular}, $\nu \perp \mu$, if there is $S \in \mathcal{A}$ with $\mu(S) = 0$ and $\nu(X \setminus S) = 0$: they live on disjoint sets.

\begin{theorem}[Lebesgue decomposition]\label{lebesgue-differentiation:theorem-lebesgue-decomposition}
Let $\mu$, $\nu$ be $\sigma$-finite measures on $(X, \mathcal{A})$. There are unique measures $\nu_a \ll \mu$ and $\nu_s \perp \mu$ with $\nu = \nu_a + \nu_s$; and $\nu_a = h\,\mu$ for a measurable $h \geq 0$.
\end{theorem}

\begin{proof}
As in Theorem~\ref{lebesgue-differentiation:theorem-radon-nikodym} we may assume $\mu$, $\nu$ finite. Let $\rho = \mu + \nu$. Since $\nu \ll \rho$, Theorem~\ref{lebesgue-differentiation:theorem-radon-nikodym} gives $g \geq 0$ with
$\nu(E) = \int_Eg\,d\rho$; from $\nu(E) \leq \rho(E)$ for all $E$, $g \leq 1$ $\rho$-almost everywhere, and we may take $0 \leq g \leq 1$. Writing $\int_Eg\,d\rho = \int_Eg\,d\mu + \int_Eg\,d\nu$ gives
\[
\int_E(1 - g)\,d\nu = \int_Eg\,d\mu \qquad (E \in \mathcal{A}).
\]
Let $S = \{g = 1\}$. For $E = S$ the left side vanishes, so $\int_Sg\,d\mu = \mu(S) = 0$. Put $\nu_s(E) = \nu(E \cap S)$ and $\nu_a(E) = \nu(E \setminus S)$; then $\nu_s \perp \mu$. If $\mu(E) = 0$, the identity gives
$\int_E(1 - g)\,d\nu = 0$, so $1 - g = 0$ $\nu$-almost everywhere on $E$, that is, $\nu(E \setminus S) = 0$: $\nu_a \ll \mu$, and $\nu_a = h\,\mu$ by Theorem~\ref{lebesgue-differentiation:theorem-radon-nikodym}. For
uniqueness, let $\nu = \nu_a^{\prime} + \nu_s^{\prime}$ be another decomposition, with $\nu_s^{\prime}$ living on a $\mu$-null set $S^{\prime}$. On subsets of $X \setminus (S \cup S^{\prime})$ both singular parts vanish,
so $\nu_a = \nu = \nu_a^{\prime}$; on subsets of the $\mu$-null set $S \cup S^{\prime}$ both absolutely continuous parts vanish. Hence $\nu_a = \nu_a^{\prime}$, and $\nu_s = \nu - \nu_a = \nu_s^{\prime}$ ($\nu$ being finite).
\end{proof}

\begin{example}\label{lebesgue-differentiation:example-lebesgue-decomposition}
With respect to Lebesgue measure $\lambda$ on $\RR$: $\nu = \delta_0 + \mathbf{1}_{[0,1]}\lambda$ has $\nu_a = \mathbf{1}_{[0,1]}\lambda$ and $\nu_s = \delta_0$, living on $\{0\}$. The measure $\mu_F$ of the
Cantor function (Lemma~\ref{lebesgue-differentiation:lemma-stieltjes}) is singular, living on the Cantor set, but has no atoms. In general the singular part splits further into a \emph{discrete} part, a countable sum of
point masses, and a \emph{singular continuous} part without atoms, like $\mu_F$: the atoms of a $\sigma$-finite measure form a countable set, and the discrete part is the restriction to it.
\end{example}

\section{Absolutely continuous functions}\label{lebesgue-differentiation:section-absolutely-continuous}

By Corollary~\ref{lebesgue-differentiation:corollary-ftc-lebesgue}, an indefinite integral $F(x) = F(a) + \int_a^xf$ of an integrable function satisfies $F' = f$ almost everywhere.
Conversely, when is a function the integral of its derivative? Continuity and differentiability almost everywhere are not enough, by the Cantor function
(Counterexample~\ref{lebesgue-differentiation:counterexample-cantor-function}). The right condition is absolute continuity.

\begin{definition}\label{lebesgue-differentiation:definition-absolutely-continuous}
A function $F \colon [a, b] \to \CC$ is \emph{absolutely continuous} if for every $\varepsilon > 0$ there is $\delta > 0$ such that $\sum_k|F(b_k) - F(a_k)| < \varepsilon$ for every finite
family of pairwise disjoint intervals $(a_k, b_k) \subseteq [a, b]$ with $\sum_k(b_k - a_k) < \delta$. For real $F$, the \emph{total variation} of $F$ on $[a, x]$ is
$V_F(x) = \sup\sum_i|F(t_i) - F(t_{i-1})|$, the supremum over the partitions $a = t_0 < \cdots < t_m = x$.
\end{definition}

\begin{lemma}\label{lebesgue-differentiation:lemma-absolutely-continuous}
\begin{enumerate}
\item For $f \in L^1([a, b])$, $F(x) = \int_{[a,x]}f\,d\lambda$ is absolutely continuous; more precisely, for every $\varepsilon > 0$ there is $\delta > 0$ with $\int_E|f| < \varepsilon$
whenever $\lambda(E) < \delta$.
\item If $F$ is real and absolutely continuous, then $V_F$ is finite, nondecreasing and absolutely continuous, and so is $V_F - F$. So $F = V_F - (V_F - F)$ is a difference of
nondecreasing absolutely continuous functions.
\end{enumerate}
\end{lemma}

\begin{proof}
(1) By monotone convergence $\int_{\{|f| > N\}}|f| \to 0$ as $N \to \infty$; choose $N$ with $\int_{\{|f| > N\}}|f| < \varepsilon/2$ and $\delta = \varepsilon/(2N)$. Then
$\int_E|f| \leq \int_{\{|f| > N\}}|f| + N\lambda(E) < \varepsilon$. For disjoint intervals of total length $< \delta$, $\sum_k|F(b_k) - F(a_k)| \leq \int_{\bigcup_k(a_k, b_k)}|f| < \varepsilon$.

(2) Take $\delta$ for $\varepsilon = 1$ and cut $[a, b]$ into $N$ intervals of length $< \delta$; refining a partition by these cut points does not decrease the sum, and the sum
over each of the $N$ intervals is $< 1$, so $V_F(b) \leq N$. For $x < y$, $V_F(y) - V_F(x)$ is the total variation of $F$ on $[x, y]$ (a partition of $[a, y]$ refined by $x$ splits
into partitions of $[a, x]$ and $[x, y]$), so $V_F$ is nondecreasing and $V_F(y) - V_F(x) \geq |F(y) - F(x)| \geq F(y) - F(x)$, which makes $V_F - F$ nondecreasing. If
$(a_k, b_k)$ are disjoint with $\sum_k(b_k - a_k) < \delta$ for the $\delta$ of $\varepsilon$, then any partitions of the intervals $[a_k, b_k]$ together give disjoint intervals of total
length $< \delta$, so $\sum_k(V_F(b_k) - V_F(a_k)) \leq \varepsilon$: $V_F$ is absolutely continuous, and so is $V_F - F$.
\end{proof}

\begin{lemma}\label{lebesgue-differentiation:lemma-stieltjes}
Let $F \colon [a, b] \to \RR$ be continuous and nondecreasing. There is a finite Borel measure $\mu_F$ on $[a, b]$ with $\mu_F([a, x]) = F(x) - F(a)$ for all $x$; so
$\mu_F((x, y]) = F(y) - F(x)$ and points have measure $0$. If $F$ is absolutely continuous, then $\mu_F$ is absolutely continuous with respect to $\lambda$.
\end{lemma}

\begin{proof}
Extend $F$ to $\RR$ by $F(x) = F(a)$ for $x < a$ and $F(x) = F(b)$ for $x > b$; the extension is continuous and nondecreasing. Let $\mu_F$ be the restriction to $[a, b]$ of its
Lebesgue--Stieltjes measure (Theorem~\ref{measure-theory:theorem-lebesgue-stieltjes}). Points have measure $0$ by part (4) of that theorem, since $F$ is continuous, and
$\mu_F([a, x]) = F(x) - F(a-) = F(x) - F(a)$.

Let $F$ be absolutely continuous and $E$ a Borel set with $\lambda(E) = 0$. Given $\varepsilon > 0$, take $\delta$ from the definition and an open $U \supseteq E$ with $\lambda(U) < \delta$
(Theorem~\ref{measure-theory:theorem-lebesgue-measure}). The open set $U \cap (a, b)$ is the disjoint union of countably many open intervals $(a_k, b_k)$, its connected components,
with $\sum_k(b_k - a_k) < \delta$. For every $N$, $\sum_{k \leq N}\mu_F((a_k, b_k)) = \sum_{k \leq N}(F(b_k) - F(a_k)) < \varepsilon$, so $\mu_F(U \cap (a, b)) \leq \varepsilon$ and
$\mu_F(E) \leq \varepsilon$, the endpoints $a$, $b$ having measure $0$. As $\varepsilon$ is arbitrary, $\mu_F(E) = 0$.
\end{proof}

\begin{theorem}[Fundamental theorem of calculus for absolutely continuous functions]\label{lebesgue-differentiation:theorem-absolutely-continuous}
Let $F \colon [a, b] \to \CC$ be absolutely continuous. Then $F$ is differentiable almost everywhere, $F' \in L^1([a, b])$, and
\[
F(x) = F(a) + \int_a^xF'(t)\,dt \qquad \text{for all } x \in [a, b].
\]
Conversely, every function of the form $F(a) + \int_a^xf$ with $f \in L^1([a, b])$ is absolutely continuous. So the absolutely continuous functions are exactly the indefinite
Lebesgue integrals.
\end{theorem}

\begin{proof}
The converse is Lemma~\ref{lebesgue-differentiation:lemma-absolutely-continuous}(1). For the direct statement, treat real and imaginary parts separately, and by
Lemma~\ref{lebesgue-differentiation:lemma-absolutely-continuous}(2) assume that $F$ is nondecreasing, the statement being linear in $F$. By
Lemma~\ref{lebesgue-differentiation:lemma-stieltjes}, $\mu_F$ is a finite measure with $\mu_F \ll \lambda$ on the Borel sets of $[a, b]$, so by the Radon--Nikodym theorem
(Theorem~\ref{lebesgue-differentiation:theorem-radon-nikodym}) there is an integrable $h \geq 0$ with $\mu_F(E) = \int_Eh\,d\lambda$. Hence $F(x) - F(a) = \mu_F([a, x]) = \int_a^xh$, and by
Corollary~\ref{lebesgue-differentiation:corollary-ftc-lebesgue}, $F$ is differentiable almost everywhere with $F' = h$.
\end{proof}

\begin{example}\label{lebesgue-differentiation:example-absolutely-continuous}
A Lipschitz function, $|F(x) - F(y)| \leq L|x - y|$, is absolutely continuous (take $\delta = \varepsilon/L$), so it is differentiable almost everywhere and is the integral of its
derivative; for instance $|x| = \int_0^x\operatorname{sgn}$. The Cantor function is continuous and nondecreasing but not absolutely continuous: it maps the Cantor set, of measure
$0$, onto $[0, 1]$, whereas the proof of Lemma~\ref{lebesgue-differentiation:lemma-stieltjes} shows that an absolutely continuous nondecreasing function gives measure $0$ to null sets, and
$\mu_F(C) = F(1) - F(0) = 1$ for the Cantor function because $F$ is constant on the complement of $C$. The function $x\sin(1/x)$ (and $0$ at $0$) on $[0, 1]$ is continuous but not
of bounded variation, hence not absolutely continuous.
\end{example}

\section{Monotone functions and bounded variation}\label{lebesgue-differentiation:section-bounded-variation}

Absolutely continuous functions are differences of nondecreasing ones
(Lemma~\ref{lebesgue-differentiation:lemma-absolutely-continuous}). Here we study the larger class of all differences of nondecreasing functions, the functions of bounded variation, without any
continuity assumption. We prove Lebesgue's theorem that they are differentiable almost everywhere, decompose them into an absolutely continuous, a singular continuous and a jump part, and characterise
the absolutely continuous ones among them by the theorem of Banach and Zarecki: a continuous function of bounded variation is absolutely continuous exactly when it maps null sets to null sets. All
functions in this section are real-valued; a complex function is treated through its real and imaginary parts.

\subsection{Functions of bounded variation}\label{lebesgue-differentiation:subsection-bv}

\begin{definition}\label{lebesgue-differentiation:definition-bounded-variation}
Let $F \colon [c, d] \to \RR$. The \emph{total variation} of $F$ on $[c, d]$ is
\[
V(F; [c, d]) = \sup\Bigl\{\sum_{i=1}^m|F(t_i) - F(t_{i-1})| : c = t_0 < t_1 < \cdots < t_m = d\Bigr\} \in [0, \infty],
\]
the supremum over all partitions of $[c, d]$. The function $F$ is \emph{of bounded variation}, $F \in BV([c, d])$, if $V(F; [c, d]) < \infty$.
\end{definition}

\begin{notation}\label{lebesgue-differentiation:notation-bounded-variation}
For $F \colon [a, b] \to \RR$ we keep the notation $V_F(x) = V(F; [a, x])$ of Definition~\ref{lebesgue-differentiation:definition-absolutely-continuous}, with $V_F(a) = 0$. One-sided limits are
written $F(x+) = \lim_{t \downarrow x}F(t)$ and $F(x-) = \lim_{t \uparrow x}F(t)$, with the conventions $F(a-) = F(a)$ and $F(b+) = F(b)$. The \emph{jump} of $F$ at $x$ is $F(x+) - F(x-)$.
\end{notation}

Adding a point to a partition does not decrease the sum, by the triangle inequality; so the supremum may be taken over partitions containing any given finite set of points, and it is the
limit of the sums along any sequence of partitions each refining the previous one and with mesh tending to $0$, provided $F$ is continuous.

\begin{example}\label{lebesgue-differentiation:example-bounded-variation}
\begin{enumerate}
\item A monotone $F$ has $V(F; [c, d]) = |F(d) - F(c)|$: all the sums telescope.
\item A Lipschitz $F$, $|F(x) - F(y)| \leq L|x - y|$, has $V(F; [c, d]) \leq L(d - c)$. In particular $C^1$ functions on $[c, d]$ are of bounded variation, by the mean value theorem.
\item The indicator $\mathbf{1}_{[1/2, 1]}$ on $[0, 1]$ has variation $1$: a sum has exactly one non-zero term. Functions of bounded variation need not be continuous.
\item $F(x) = x\sin(1/x)$, $F(0) = 0$, is continuous on $[0, 1]$ but not of bounded variation. At $t_k = 1/((k + \frac12)\pi)$ it takes the values $\pm t_k$ with alternating signs, so
$|F(t_k) - F(t_{k+1})| = t_k + t_{k+1} \geq t_k$, and $\sum_k t_k = \infty$.
\item $F(x) = x^2\sin(1/x^2)$, $F(0) = 0$, is differentiable at every point of $[0, 1]$, but the same argument with $t_k = ((k + \frac12)\pi)^{-1/2}$ gives the terms $t_k^2 + t_{k+1}^2 \geq t_k^2$,
whose sum diverges. In contrast, $x^2\sin(1/x)$ has $|F^{\prime}| \leq 3$ on $[0, 1]$, so it is Lipschitz and of bounded variation.
\end{enumerate}
\end{example}

\begin{proposition}\label{lebesgue-differentiation:proposition-variation}
Let $F, G \colon [a, b] \to \RR$.
\begin{enumerate}
\item For $a \leq c \leq d \leq e \leq b$, $V(F; [c, e]) = V(F; [c, d]) + V(F; [d, e])$, and $|F(e) - F(c)| \leq V(F; [c, e])$.
\item Functions of bounded variation are bounded: $|F(x)| \leq |F(a)| + V(F; [a, b])$.
\item $V(F + G; [a, b]) \leq V(F; [a, b]) + V(G; [a, b])$ and $V(tF; [a, b]) = |t|V(F; [a, b])$; so $BV([a, b])$ is a vector space. It is closed under
products:
\[
V(FG; [a, b]) \leq \|F\|_\infty V(G; [a, b]) + \|G\|_\infty V(F; [a, b]).
\]
\end{enumerate}
\end{proposition}

\begin{proof}
(1) A partition of $[c, e]$ refined by $d$ splits into partitions of $[c, d]$ and $[d, e]$, and conversely two such partitions together form a partition of $[c, e]$ containing $d$; since adding $d$ does not
decrease the sum, the suprema agree. The inequality is the trivial partition. (2) follows from (1): $|F(x) - F(a)| \leq V(F; [a, x]) \leq V(F; [a, b])$. (3) The first two statements hold for each partition sum.
For the product, $|F(t_i)G(t_i) - F(t_{i-1})G(t_{i-1})| \leq |F(t_i)||G(t_i) - G(t_{i-1})| + |G(t_{i-1})||F(t_i) - F(t_{i-1})|$.
\end{proof}

\begin{theorem}[Jordan decomposition]\label{lebesgue-differentiation:theorem-jordan}
A function $F \colon [a, b] \to \RR$ is of bounded variation if and only if it is a difference of two nondecreasing functions. More precisely, if $F \in BV([a, b])$, the \emph{positive} and
\emph{negative variations}
\[
P_F = \tfrac12(V_F + F - F(a)), \qquad N_F = \tfrac12(V_F - F + F(a))
\]
are nondecreasing, vanish at $a$, and satisfy $F = F(a) + P_F - N_F$ and $V_F = P_F + N_F$. They are minimal: if $F = G - H$ with $G$, $H$ nondecreasing, then
$P_F(y) - P_F(x) \leq G(y) - G(x)$ and $N_F(y) - N_F(x) \leq H(y) - H(x)$ for $x < y$.
\end{theorem}

\begin{proof}
A difference of nondecreasing functions is of bounded variation by Example~\ref{lebesgue-differentiation:example-bounded-variation}(1) and Proposition~\ref{lebesgue-differentiation:proposition-variation}(3).
Conversely, let $F \in BV([a, b])$ and $x < y$. By Proposition~\ref{lebesgue-differentiation:proposition-variation}(1), $V_F(y) - V_F(x) = V(F; [x, y]) \geq |F(y) - F(x)|$, so both
$V_F(y) - V_F(x) \pm (F(y) - F(x))$ are $\geq 0$: $P_F$ and $N_F$ are nondecreasing. The identities are immediate. For minimality, write $\Delta_iF = F(t_i) - F(t_{i-1})$ for a partition of $[x, y]$. Then
$|\Delta_iF| \pm \Delta_iF = 2(\Delta_iF)^{\pm} \leq 2\Delta_iG$ for $+$ and $\leq 2\Delta_iH$ for $-$, because $\Delta_iF = \Delta_iG - \Delta_iH$ with both increments $\geq 0$. Summing,
$\sum_i|\Delta_iF| + (F(y) - F(x)) \leq 2(G(y) - G(x))$, and taking the supremum over partitions, $2(P_F(y) - P_F(x)) \leq 2(G(y) - G(x))$; similarly for $N_F$.
\end{proof}

The Jordan decomposition of a function is the analogue of the Jordan decomposition of a signed measure into positive and negative parts (Chapter~\ref{signed-measures}), and the two correspond under
$F \mapsto \mu_F$; the minimality statement is the analogue of the minimality of $\mu^+$ and $\mu^-$.

\subsection{Lebesgue's theorem on monotone functions}\label{lebesgue-differentiation:subsection-monotone}

\begin{lemma}\label{lebesgue-differentiation:lemma-monotone-jumps}
Let $F \colon [a, b] \to \RR$ be nondecreasing. Then $F(x-)$ and $F(x+)$ exist for every $x$, with $F(x-) \leq F(x) \leq F(x+)$; $F$ is continuous at $x$ if and only if its jump at $x$ is $0$; and for
finitely many points $a \leq s_1 < \cdots < s_m \leq b$ the jumps satisfy $\sum_j(F(s_j+) - F(s_j-)) \leq F(b) - F(a)$. Consequently $F$ has at most countably many points of discontinuity, and the sum of all
its jumps is at most $F(b) - F(a)$.
\end{lemma}

\begin{proof}
$F(x-) = \sup_{t < x}F(t)$ and $F(x+) = \inf_{t > x}F(t)$ by monotonicity, and both are squeezed around $F(x)$. Choose points $a \leq u_0 \leq s_1 \leq u_1 \leq s_2 \leq \cdots \leq s_m \leq u_m \leq b$ with
$u_{j-1} < s_j < u_j$ whenever $s_j$ is an interior point (and $u_{j-1} = s_j$ or $u_j = s_j$ at an endpoint). Then $F(s_j+) - F(s_j-) \leq F(u_j) - F(u_{j-1})$, and the right sides telescope to at most
$F(b) - F(a)$. So at most $k(F(b) - F(a))$ points have a jump larger than $1/k$, and the set of discontinuities is a countable union of finite sets.
\end{proof}

\begin{definition}\label{lebesgue-differentiation:definition-dini}
For $F \colon [a, b] \to \RR$ and $x \in (a, b)$ the four \emph{Dini derivatives} are
\[\begin{aligned}
D^+F(x) &= \limsup_{h \downarrow 0}\frac{F(x + h) - F(x)}{h}, & D_+F(x) &= \liminf_{h \downarrow 0}\frac{F(x + h) - F(x)}{h}, \\
D^-F(x) &= \limsup_{h \downarrow 0}\frac{F(x) - F(x - h)}{h}, & D_-F(x) &= \liminf_{h \downarrow 0}\frac{F(x) - F(x - h)}{h},
\end{aligned}\]
with values in $[-\infty, \infty]$. Always $D_+F \leq D^+F$ and $D_-F \leq D^-F$, and $F$ is differentiable at $x$ if and only if the four are equal and finite.
\end{definition}

For $F(x) = |x|$ at $0$: $D^+F = D_+F = 1$ and $D^-F = D_-F = -1$. For $F(x) = x\sin(1/x)$ at $0$: $D^+F = D^-F = 1$ and $D_+F = D_-F = -1$.

\begin{theorem}[Lebesgue]\label{lebesgue-differentiation:theorem-lebesgue-monotone}
Let $F \colon [a, b] \to \RR$ be nondecreasing. Then $F$ is differentiable almost everywhere, $F^{\prime}$ is integrable, and
\[
\int_a^bF^{\prime}\,d\lambda \leq F(b) - F(a).
\]
\end{theorem}

\begin{proof}
\emph{Step 1: $D^+F \leq D_-F$ almost everywhere.} For rationals $0 \leq u < v$ let $E = \{x \in (a, b) : D_-F(x) < u < v < D^+F(x)\}$ and $s = \lambda^*(E)$. Let $\varepsilon > 0$ and $U \supseteq E$ open with
$\lambda(U) < s + \varepsilon$. For $x \in E$ there are arbitrarily small $h > 0$ with $[x - h, x] \subseteq U$ and $F(x) - F(x - h) < uh$; these intervals form a fine cover of $E$, so by
Corollary~\ref{lebesgue-differentiation:corollary-vitali-fine} there are disjoint ones $I_1 = [x_1 - h_1, x_1], \ldots, I_m = [x_m - h_m, x_m]$ with $\lambda^*(E \setminus \bigcup_iI_i) < \varepsilon$. Then
\[
\sum_i(F(x_i) - F(x_i - h_i)) < u\sum_ih_i \leq u\lambda(U) < u(s + \varepsilon).
\]
Let $A = E \cap \bigcup_i\operatorname{int}I_i$; it differs from $E \cap \bigcup_iI_i$ by finitely many points, so $\lambda^*(A) > s - \varepsilon$. For $y \in A$ there are arbitrarily small $k > 0$ such that $[y, y + k]$ lies
in the interior of the interval $I_i$ containing $y$, and $F(y + k) - F(y) > vk$. By the corollary again there are disjoint such intervals $J_1 = [y_1, y_1 + k_1], \ldots, J_p = [y_p, y_p + k_p]$ with
$\lambda^*(A \setminus \bigcup_jJ_j) < \varepsilon$, so $\sum_jk_j > s - 2\varepsilon$. Each $J_j$ lies in some $I_i$, and since $F$ is nondecreasing, the increments of $F$ over the disjoint $J_j$ inside $I_i$ add
up to at most the increment over $I_i$. Hence
\[
v(s - 2\varepsilon) < \sum_j(F(y_j + k_j) - F(y_j)) \leq \sum_i(F(x_i) - F(x_i - h_i)) < u(s + \varepsilon).
\]
Letting $\varepsilon \to 0$, $vs \leq us$, so $s = 0$. The set where $D^+F > D_-F$ is the union of the sets $E$ over all rational pairs (note $D_-F \geq 0$), so it is null.

\emph{Step 2: the derivative exists in $[0, \infty]$ almost everywhere.} The function $G(x) = -F(-x)$ on $[-b, -a]$ is nondecreasing, and $D^+G(-x) = D^-F(x)$, $D_-G(-x) = D_+F(x)$. Step 1 for $G$ gives
$D^-F \leq D_+F$ almost everywhere. So almost everywhere
\[
D^+F \leq D_-F \leq D^-F \leq D_+F \leq D^+F,
\]
all four Dini derivatives agree, and the limit $F^{\prime}(x) = \lim_{h \to 0}(F(x + h) - F(x))/h$ exists in $[0, \infty]$.

\emph{Step 3: integrability.} Extend $F$ by $F(x) = F(b)$ for $x > b$, and let $g_k(x) = k(F(x + 1/k) - F(x))$. The $g_k$ are nonnegative and Borel (monotone functions are Borel, their superlevel sets being
intervals), and $g_k \to F^{\prime}$ almost everywhere on $[a, b]$. By Fatou's lemma and translation invariance,
\[\begin{aligned}
\int_a^bF^{\prime} &\leq \liminf_k\int_a^bg_k = \liminf_k\Bigl(k\int_b^{b + 1/k}F - k\int_a^{a + 1/k}F\Bigr) \\
&= \liminf_k\Bigl(F(b) - k\int_a^{a + 1/k}F\Bigr) \leq F(b) - F(a),
\end{aligned}\]
since $F \geq F(a)$ on $[a, a + 1/k]$. In particular $F^{\prime} < \infty$ almost everywhere, so $F$ is differentiable almost everywhere.
\end{proof}

\begin{corollary}\label{lebesgue-differentiation:corollary-bv-differentiable}
Every $F \in BV([a, b])$ is differentiable almost everywhere, $F^{\prime} \in L^1([a, b])$, and $\int_a^b|F^{\prime}|\,d\lambda \leq V(F; [a, b])$.
\end{corollary}

\begin{proof}
By Theorem~\ref{lebesgue-differentiation:theorem-jordan}, $F = F(a) + P_F - N_F$; both are differentiable almost everywhere by Theorem~\ref{lebesgue-differentiation:theorem-lebesgue-monotone}, and where both are,
$|F^{\prime}| \leq P_F^{\prime} + N_F^{\prime}$. Integrating, $\int|F^{\prime}| \leq P_F(b) + N_F(b) = V_F(b)$.
\end{proof}

A series of nondecreasing functions may be differentiated term by term, a fact that also follows from Lebesgue's theorem.

\begin{theorem}[Fubini's theorem on series of monotone functions]\label{lebesgue-differentiation:theorem-fubini-series}
Let $F_n \colon [a, b] \to \RR$ be nondecreasing functions such that $F(x) = \sum_nF_n(x)$ converges for every $x \in [a, b]$. Then $F^{\prime} = \sum_nF_n^{\prime}$ almost everywhere.
\end{theorem}

\begin{proof}
\emph{Step 1: normalisation and a first inequality.} Replacing $F_n$ by $F_n - F_n(a)$, which changes neither the convergence nor the derivatives, we may assume $F_n(a) = 0$, so $F_n \geq 0$. The partial
sums $S_N = \sum_{n \leq N}F_n$ and the remainders $R_N = F - S_N = \sum_{n > N}F_n$ are nondecreasing, and $0 \leq R_N \leq R_N(b) \to 0$. Outside a null set $Z$, the countably many functions $F$, $F_n$,
$R_N$ are all differentiable by Theorem~\ref{lebesgue-differentiation:theorem-lebesgue-monotone}, and there $F_{N+1}^{\prime} \geq 0$ and $R_N^{\prime} \geq 0$ give
$S_N^{\prime} \leq S_{N+1}^{\prime} \leq F^{\prime}$. So $\sum_nF_n^{\prime}$ converges outside $Z$, with sum at most $F^{\prime}$.

\emph{Step 2: a subsequence of remainders.} Choose $N_1 < N_2 < \cdots$ with $R_{N_k}(b) < 2^{-k}$. The series $G = \sum_kR_{N_k}$ of nondecreasing nonnegative functions converges everywhere, as
$R_{N_k}(x) \leq 2^{-k}$. By Step 1 applied to it, $\sum_kR_{N_k}^{\prime} \leq G^{\prime} < \infty$ almost everywhere, so $R_{N_k}^{\prime} \to 0$, that is, $S_{N_k}^{\prime} \to F^{\prime}$, almost
everywhere. Since $S_N^{\prime}$ increases with $N$, the whole sequence converges to $F^{\prime}$ almost everywhere.
\end{proof}

The inequality in Theorem~\ref{lebesgue-differentiation:theorem-lebesgue-monotone} is strict for the Cantor function (Counterexample~\ref{lebesgue-differentiation:counterexample-cantor-function}), and for
every function with a jump. The next theorem shows that these are the only two obstructions.

\begin{definition}\label{lebesgue-differentiation:definition-jump-function}
Let $F \colon [a, b] \to \RR$ be nondecreasing. Its \emph{jump function} is
\[
J_F(x) = \sum_{a \leq t < x}(F(t+) - F(t)) + \sum_{a < t \leq x}(F(t) - F(t-)),
\]
the sums running over the countably many discontinuities (Lemma~\ref{lebesgue-differentiation:lemma-monotone-jumps}). A nondecreasing $F$ is \emph{singular} if $F^{\prime} = 0$ almost everywhere.
\end{definition}

\begin{theorem}[Lebesgue decomposition of a monotone function]\label{lebesgue-differentiation:theorem-monotone-decomposition}
Let $F \colon [a, b] \to \RR$ be nondecreasing, and put $F_{ac}(x) = F(a) + \int_a^xF^{\prime}\,d\lambda$. Then
\[
F = F_{ac} + F_{sc} + J_F,
\]
where $F_{ac}$ is absolutely continuous and nondecreasing, $J_F$ is nondecreasing with $J_F^{\prime} = 0$ almost everywhere, and $F_{sc}$ is continuous, nondecreasing and singular. The decomposition is unique:
if $F = A + S + J_F$ with $A$ absolutely continuous, $A(a) = F(a)$, and $S$ continuous and singular, then $A = F_{ac}$ and $S = F_{sc}$.
\end{theorem}

\begin{proof}
\emph{Step 1: $F - F_{ac}$ is nondecreasing and singular.} For $x < y$, Theorem~\ref{lebesgue-differentiation:theorem-lebesgue-monotone} on $[x, y]$ gives $\int_x^yF^{\prime} \leq F(y) - F(x)$, which says that
$F - F_{ac}$ is nondecreasing. By Corollary~\ref{lebesgue-differentiation:corollary-ftc-lebesgue}, $F_{ac}^{\prime} = F^{\prime}$ almost everywhere, so $(F - F_{ac})^{\prime} = 0$ almost everywhere. $F_{ac}$ is absolutely
continuous by Lemma~\ref{lebesgue-differentiation:lemma-absolutely-continuous}(1), and nondecreasing as $F^{\prime} \geq 0$.

\emph{Step 2: $J_F$ and $F - J_F$ are nondecreasing.} For $x < y$,
\[
J_F(y) - J_F(x) = \sum_{x \leq t < y}(F(t+) - F(t)) + \sum_{x < t \leq y}(F(t) - F(t-)) \geq 0 .
\]
Any finitely many terms of this sum are bounded by $F(y) - F(x)$: choosing intermediate points as in the proof of Lemma~\ref{lebesgue-differentiation:lemma-monotone-jumps}, they are bounded by increments of $F$
over disjoint subintervals of $[x, y]$. So $J_F(y) - J_F(x) \leq F(y) - F(x)$.

\emph{Step 3: $F - J_F$ is continuous.} Let $x < y$. Since $F(y-) \geq F(x+)$, the term $F(y) - F(y-)$ is at most $F(y) - F(x+)$, which tends to $0$ as $y \downarrow x$; and the sum of the jumps at points
$t \in (x, y)$ tends to $0$ as $y \downarrow x$, being the tail of a convergent series of nonnegative terms. So $J_F(x+) - J_F(x) = F(x+) - F(x)$, and similarly $J_F(x) - J_F(x-) = F(x) - F(x-)$. Hence
$F - J_F$ has no jumps, and it is continuous by Lemma~\ref{lebesgue-differentiation:lemma-monotone-jumps}.

\emph{Step 4: conclusion.} Put $F_{sc} = F - F_{ac} - J_F$, continuous as $F - J_F$ and $F_{ac}$ are. The function $F - F_{ac}$ is nondecreasing and has the same jumps as $F$, $F_{ac}$ being continuous;
applying Step 2 to it, $F_{sc} = (F - F_{ac}) - J_{F - F_{ac}}$ is nondecreasing. Now $F_{sc}$ and $J_F$ are nondecreasing, so their derivatives, which exist almost everywhere, are $\geq 0$; they add up to
$(F - F_{ac})^{\prime} = 0$ almost everywhere by Step 1. So both vanish almost everywhere.

\emph{Step 5: uniqueness.} If $F = A + S + J_F$ as stated, then $A^{\prime} = F^{\prime} - S^{\prime} - J_F^{\prime} = F^{\prime}$ almost everywhere, and by
Theorem~\ref{lebesgue-differentiation:theorem-absolutely-continuous}, $A(x) = F(a) + \int_a^xA^{\prime} = F_{ac}(x)$. Then $S = F - F_{ac} - J_F = F_{sc}$.
\end{proof}

\begin{example}\label{lebesgue-differentiation:example-monotone-decomposition}
Let $c$ be the Cantor function on $[0, 1]$ and $F(x) = x + c(x) + \mathbf{1}_{[1/2, 1]}(x)$. Then $F^{\prime} = 1$ almost everywhere, $F_{ac}(x) = x$, $J_F = \mathbf{1}_{[1/2, 1]}$ (the only jump, of size $1$,
is at $\frac12$, where $F(\frac12-) = F(\frac12) - 1$), and $F_{sc} = c$. Under $F \mapsto \mu_F$ (Theorem~\ref{measure-theory:theorem-lebesgue-stieltjes}, after making $F$ right-continuous) the three parts
correspond to the absolutely continuous part $\lambda$, the singular continuous part $\mu_c$ and the discrete part $\delta_{1/2}$ of the measure, as in
Example~\ref{lebesgue-differentiation:example-lebesgue-decomposition}.
\end{example}

\subsection{The theorem of Banach and Zarecki}\label{lebesgue-differentiation:subsection-banach-zarecki}

What distinguishes the Cantor function from an absolutely continuous function? It is continuous and of bounded variation, but it maps the Cantor set, a null set, onto $[0, 1]$. This turns out to be the
only possible defect.

\begin{definition}\label{lebesgue-differentiation:definition-luzin-n}
A function $F \colon [a, b] \to \RR$ has \emph{Luzin's property (N)} if $\lambda(F(Z)) = 0$ for every null set $Z \subseteq [a, b]$ (the image $F(Z)$ is then measurable, being a null set).
\end{definition}

\begin{lemma}\label{lebesgue-differentiation:lemma-image-measure-bound}
Let $F \colon [a, b] \to \RR$ and $E \subseteq [a, b]$.
\begin{enumerate}
\item If $F$ is differentiable at every point of $E$ with $|F^{\prime}| \leq M$ on $E$, then $\lambda^*(F(E)) \leq M\lambda^*(E)$.
\item If $F$ is continuous, $E$ is measurable and $F$ is differentiable at every point of $E$, then $F^{\prime}$ is measurable on $E$ and $\lambda^*(F(E)) \leq \int_E|F^{\prime}|\,d\lambda$.
\end{enumerate}
\end{lemma}

\begin{proof}
(1) Let $\varepsilon > 0$, and for $k \geq 1$ let $E_k$ be the set of $x \in E$ with $|F(y) - F(x)| \leq (M + \varepsilon)|y - x|$ for all $y \in [a, b]$ with $|y - x| < 1/k$. The $E_k$ increase, with union $E$ by
the definition of the derivative. Cover $E_k$ by countably many intervals $I_j$ of length $< 1/k$ with $\sum_j\lambda(I_j) \leq \lambda^*(E_k) + \varepsilon$. Two points $x, y$ of $E_k \cap I_j$ satisfy
$|F(x) - F(y)| \leq (M + \varepsilon)\lambda(I_j)$, so $F(E_k \cap I_j)$ lies in an interval of length $(M + \varepsilon)\lambda(I_j)$, and
$\lambda^*(F(E_k)) \leq (M + \varepsilon)(\lambda^*(E) + \varepsilon)$. It remains to pass to the limit in $k$. Choose Borel sets $H_k \supseteq F(E_k)$ with $\lambda(H_k) = \lambda^*(F(E_k))$ (intersections of
open sets), and replace $H_k$ by $\bigcap_{j \geq k}H_j$, which still contains $F(E_k)$ since the $F(E_j)$ increase; now the $H_k$ increase, and their union contains $F(E)$, so
$\lambda^*(F(E)) \leq \lim_k\lambda(H_k) \leq (M + \varepsilon)(\lambda^*(E) + \varepsilon)$. Let $\varepsilon \to 0$.

(2) Extend $F$ by $F(x) = F(b)$ for $x > b$. The functions $g_k(x) = k(F(x + 1/k) - F(x))$ are continuous, and at every point of $E \setminus \{b\}$ they converge to $F^{\prime}$; so $F^{\prime}$ is
measurable on $E$. For $\eta > 0$ and $m \geq 1$ let $E_m = \{x \in E : (m - 1)\eta \leq |F^{\prime}(x)| < m\eta\}$, a measurable set. By (1), $\lambda^*(F(E_m)) \leq m\eta\lambda(E_m)$, and so
\[
\lambda^*(F(E)) \leq \sum_m m\eta\lambda(E_m) \leq \sum_m\int_{E_m}(|F^{\prime}| + \eta)\,d\lambda = \int_E|F^{\prime}|\,d\lambda + \eta\lambda(E);
\]
let $\eta \to 0$.
\end{proof}
\begin{theorem}[Banach--Zarecki]\label{lebesgue-differentiation:theorem-banach-zarecki}
A function $F \colon [a, b] \to \RR$ is absolutely continuous if and only if it is continuous, of bounded variation, and has Luzin's property (N).
\end{theorem}

\begin{proof}
\emph{Absolutely continuous functions have the three properties.} Continuity is the definition applied to a single interval, and bounded variation is
Lemma~\ref{lebesgue-differentiation:lemma-absolutely-continuous}(2). Let $Z$ be a null set, $\varepsilon > 0$, and $\delta$ as in the definition of absolute continuity. Choose an open $U$ with
$Z \cap (a, b) \subseteq U \subseteq (a, b)$ and $\lambda(U) < \delta$, and write $U$ as the disjoint union of its components $(a_k, b_k)$. On $[a_k, b_k]$ the continuous $F$ attains its minimum and maximum at points
$s_k$, $t_k$, so $F([a_k, b_k])$ is an interval of length $|F(t_k) - F(s_k)|$. The open intervals between $s_k$ and $t_k$ lie in the disjoint $(a_k, b_k)$, and their total length is $< \delta$; so for every $m$,
$\sum_{k \leq m}|F(t_k) - F(s_k)| < \varepsilon$. Hence
\[
\lambda^*(F(Z)) \leq \lambda(F(\{a, b\})) + \sum_k\lambda(F([a_k, b_k])) \leq \varepsilon,
\]
and $F(Z)$ is null.

\emph{The three properties imply absolute continuity.} By Corollary~\ref{lebesgue-differentiation:corollary-bv-differentiable}, $F$ is differentiable outside a null set $Z$, and $F^{\prime} \in L^1([a, b])$. Let
$x < y$ in $[a, b]$. By the intermediate value theorem $F([x, y])$ contains the interval between $F(x)$ and $F(y)$, so by Lemma~\ref{lebesgue-differentiation:lemma-image-measure-bound}(2), applied to the measurable
set $[x, y] \setminus Z$, and property (N),
\[
|F(y) - F(x)| \leq \lambda^*(F([x, y] \setminus Z)) + \lambda(F(Z)) \leq \int_x^y|F^{\prime}|\,d\lambda .
\]
Given $\varepsilon > 0$, Lemma~\ref{lebesgue-differentiation:lemma-absolutely-continuous}(1) gives $\delta > 0$ with $\int_E|F^{\prime}| < \varepsilon$ whenever $\lambda(E) < \delta$; then for disjoint intervals $(a_k, b_k)$
of total length $< \delta$, $\sum_k|F(b_k) - F(a_k)| \leq \int_{\bigcup_k(a_k, b_k)}|F^{\prime}| < \varepsilon$.
\end{proof}

\begin{counterexample}\label{lebesgue-differentiation:counterexample-banach-zarecki}
None of the three conditions of Theorem~\ref{lebesgue-differentiation:theorem-banach-zarecki} follows from the other two.
\begin{enumerate}
\item The Cantor function is continuous and of bounded variation (it is monotone), but it maps the Cantor set, a null set, onto $[0, 1]$.
\item $F(x) = x\sin(1/x)$, $F(0) = 0$, is continuous and not of bounded variation (Example~\ref{lebesgue-differentiation:example-bounded-variation}(4)), but it has property (N). On $[1/k, 1]$ it is $C^1$, so
$|F^{\prime}| \leq M_k$ there and Lemma~\ref{lebesgue-differentiation:lemma-image-measure-bound}(1) gives $\lambda(F(Z \cap [1/k, 1])) = 0$ for a null set $Z$; and
$F(Z) \subseteq \{0\} \cup \bigcup_kF(Z \cap [1/k, 1])$.
\item $\mathbf{1}_{[1/2, 1]}$ is of bounded variation and has property (N), as all its images are finite, but it is not continuous.
\end{enumerate}
\end{counterexample}

\begin{remark}\label{lebesgue-differentiation:remark-luzin-measurable}
A continuous $F \colon [a, b] \to \RR$ with property (N) maps measurable sets to measurable sets. Indeed, by inner regularity a measurable $E \subseteq [a, b]$ is the union of countably many compact sets $K_i$ and
a null set $Z$; the images $F(K_i)$ are compact and $F(Z)$ is null. Conversely, if a continuous $F$ maps measurable sets to measurable sets, it has property (N): if $\lambda(F(Z)) > 0$ for a null set $Z$, then
either $B = F(Z)$ is not measurable, or it is measurable of positive measure and contains a non-measurable set $B$ by Theorem~\ref{measure-theory:theorem-bernstein}; in both cases
$F^{-1}(B) \cap Z$ is a null set, hence measurable, with image $B$. The Cantor
function therefore maps some null set onto a non-measurable set.
\end{remark}

The theorem gives two useful consequences: a formula for the variation, and a fundamental theorem of calculus for functions that are differentiable everywhere.

\begin{corollary}\label{lebesgue-differentiation:corollary-variation-integral}
If $F \colon [a, b] \to \RR$ is absolutely continuous, then $V_F(x) = \int_a^x|F^{\prime}|\,d\lambda$ for all $x \in [a, b]$. In particular $V(F; [a, b]) = \int_a^b|F^{\prime}|$ for $F \in C^1([a, b])$.
\end{corollary}

\begin{proof}
For a partition of $[a, x]$, $\sum_i|F(t_i) - F(t_{i-1})| = \sum_i|\int_{t_{i-1}}^{t_i}F^{\prime}| \leq \int_a^x|F^{\prime}|$ by Theorem~\ref{lebesgue-differentiation:theorem-absolutely-continuous}; so
$V_F(x) \leq \int_a^x|F^{\prime}|$. The reverse inequality is Corollary~\ref{lebesgue-differentiation:corollary-bv-differentiable} on $[a, x]$. A $C^1$ function is Lipschitz, hence absolutely continuous.
\end{proof}

\begin{corollary}\label{lebesgue-differentiation:corollary-everywhere-differentiable}
Let $F \colon [a, b] \to \RR$ be continuous, differentiable at all but countably many points of $[a, b]$, and suppose $F^{\prime} \in L^1([a, b])$. Then $F$ is absolutely continuous, and
$F(x) = F(a) + \int_a^xF^{\prime}\,d\lambda$ for all $x$.
\end{corollary}

\begin{proof}
Let $D$ be the countable set where $F$ is not differentiable. \emph{Property (N).} For a null set $Z$ and $M \geq 1$, let $Z_M = \{x \in Z \setminus D : |F^{\prime}(x)| \leq M\}$. By
Lemma~\ref{lebesgue-differentiation:lemma-image-measure-bound}(1), $\lambda^*(F(Z_M)) \leq M\lambda^*(Z_M) = 0$; and $F(Z \cap D)$ is countable. So $F(Z) = F(Z \cap D) \cup \bigcup_MF(Z_M)$ is null.
\emph{Bounded variation.} For $x < y$, as in the proof of the theorem, $|F(y) - F(x)| \leq \lambda^*(F([x, y] \setminus D)) + \lambda(F(D)) \leq \int_x^y|F^{\prime}|$ by
Lemma~\ref{lebesgue-differentiation:lemma-image-measure-bound}(2), so $V(F; [a, b]) \leq \int_a^b|F^{\prime}| < \infty$. By Theorem~\ref{lebesgue-differentiation:theorem-banach-zarecki}, $F$ is absolutely continuous, and
Theorem~\ref{lebesgue-differentiation:theorem-absolutely-continuous} gives the formula.
\end{proof}

The hypotheses are sharp. Differentiability almost everywhere is not enough: the Cantor function is differentiable outside the Cantor set, with $F^{\prime} = 0$. Integrability of $F^{\prime}$ is needed:
$x^2\sin(1/x^2)$ is differentiable everywhere, but it is not of bounded variation (Example~\ref{lebesgue-differentiation:example-bounded-variation}(5)), and its derivative is not integrable.

\section{Exercises}\label{lebesgue-differentiation:section-exercises}

\begin{exercise}\label{lebesgue-differentiation:exercise-no-half-density}
Show that there is no Lebesgue measurable $A \subseteq \RR$ with $\lambda(A \cap I) = \frac12\lambda(I)$ for every bounded interval $I$.
\end{exercise}

\begin{solution}
For such $A$, $\lambda(A \cap B(x, r))/\lambda(B(x, r)) = \frac12$ for all $x$ and $r$. By Corollary~\ref{lebesgue-differentiation:corollary-density} this ratio tends to $1$ at
almost every point of $A$ and to $0$ at almost every point outside $A$; so $A$ and $\RR \setminus A$ are both null, which is impossible.
\end{solution}

\begin{exercise}\label{lebesgue-differentiation:exercise-maximal-not-integrable}
Show that $Mf \notin L^1(\RR^n)$ for every $f \in L^1(\RR^n)$ that is not zero almost everywhere; so the maximal inequality cannot be improved to
$\|Mf\|_1 \leq C\|f\|_1$.
\end{exercise}

\begin{solution}
Choose $R$ with $a = \int_{B(0, R)}|f| > 0$. For $\|x\| \geq R$, the ball $B(x, 2\|x\|)$ contains $B(0, R)$, so
$Mf(x) \geq a/\lambda(B(x, 2\|x\|)) = a/(2^n\|x\|^n\lambda(B(0, 1)))$. The function $\|x\|^{-n}$ is not integrable on $\{\|x\| \geq R\}$: the shells
$2^kR \leq \|x\| < 2^{k+1}R$ have measure $(2^n - 1)2^{kn}R^n\lambda(B(0, 1))$, and on each of them $\|x\|^{-n} \geq 2^{-(k+1)n}R^{-n}$, so each contributes at least the
positive constant $(1 - 2^{-n})\lambda(B(0, 1))$ to the integral.
\end{solution}

\begin{exercise}\label{lebesgue-differentiation:exercise-variation-continuous}
Let $F \in BV([a, b])$ and $x \in [a, b]$. Show that $F$ is continuous at $x$ if and only if $V_F$ is continuous at $x$. Deduce that a continuous function of bounded variation is the difference of two continuous
nondecreasing functions.
\end{exercise}

\begin{solution}
If $V_F$ is continuous at $x$, so is $F$, because $|F(y) - F(x)| \leq |V_F(y) - V_F(x)|$ by Proposition~\ref{lebesgue-differentiation:proposition-variation}(1). Conversely, let $F$ be continuous at $x < b$ and
$\varepsilon > 0$. Choose a partition $x = t_0 < t_1 < \cdots < t_m = b$ with $\sum_i|F(t_i) - F(t_{i-1})| > V(F; [x, b]) - \varepsilon$; adding a point close to $x$ does not decrease the sum, so by continuity at $x$
we may assume $|F(t_1) - F(x)| < \varepsilon$. Then $V(F; [x, b]) < 2\varepsilon + \sum_{i \geq 2}|F(t_i) - F(t_{i-1})| \leq 2\varepsilon + V(F; [t_1, b])$, so $V_F(t_1) - V_F(x) = V(F; [x, t_1]) < 2\varepsilon$, and the same
holds for $t_1$ replaced by any $y \in (x, t_1]$, as $V_F$ is nondecreasing. So $V_F$ is right-continuous at $x$; left continuity is proved in the same way with partitions of $[a, x]$. For the last statement,
$P_F$ and $N_F$ of Theorem~\ref{lebesgue-differentiation:theorem-jordan} are continuous when $F$ and $V_F$ are.
\end{solution}

\begin{exercise}\label{lebesgue-differentiation:exercise-integration-by-parts}
Let $F, G \colon [a, b] \to \RR$ be absolutely continuous. Show that $FG$ is absolutely continuous and
\[
\int_a^bFG^{\prime}\,d\lambda = F(b)G(b) - F(a)G(a) - \int_a^bF^{\prime}G\,d\lambda .
\]
\end{exercise}

\begin{solution}
$F$ and $G$ are continuous, hence bounded. For disjoint intervals $(a_k, b_k)$,
\[\begin{aligned}
\sum_k|F(b_k)G(b_k) - F(a_k)G(a_k)| &\leq \|F\|_\infty\sum_k|G(b_k) - G(a_k)| \\
&\quad + \|G\|_\infty\sum_k|F(b_k) - F(a_k)|,
\end{aligned}\]
which is small when the total length is small; so $FG$ is absolutely continuous. At every point
where both $F$ and $G$ are differentiable, which is almost every point, $(FG)^{\prime} = F^{\prime}G + FG^{\prime}$, and $F^{\prime}G$, $FG^{\prime}$ are integrable (products of integrable and bounded
measurable functions). Theorem~\ref{lebesgue-differentiation:theorem-absolutely-continuous} applied to $FG$ gives $F(b)G(b) - F(a)G(a) = \int_a^b(F^{\prime}G + FG^{\prime})$.
\end{solution}

\begin{exercise}\label{lebesgue-differentiation:exercise-composition-ac}
Let $F(x) = \sqrt{x}$ and $G(x) = x^2\sin^2(1/x)$, $G(0) = 0$, on $[0, 1]$. Show that $F$ and $G$ are absolutely continuous and $G([0, 1]) \subseteq [0, 1]$, but $F \circ G$ is not of bounded variation. So a
composition of absolutely continuous functions need not be absolutely continuous.
\end{exercise}

\begin{solution}
$G$ is differentiable on $[0, 1]$ with $G^{\prime}(x) = 2x\sin^2(1/x) - 2\sin(1/x)\cos(1/x)$ for $x > 0$ and $G^{\prime}(0) = 0$, so $|G^{\prime}| \leq 4$, $G$ is Lipschitz and hence absolutely continuous, and
$0 \leq G \leq 1$. $F(x) = \int_0^x\frac{dt}{2\sqrt t}$ with an integrable integrand, so $F$ is absolutely continuous by Theorem~\ref{lebesgue-differentiation:theorem-absolutely-continuous}. Now
$F(G(x)) = x|\sin(1/x)|$. With $s_k = 1/((k + \frac12)\pi)$ and $u_k = 1/(k\pi)$ we have $u_{k+1} < s_k < u_k$, $F(G(s_k)) = s_k$ and $F(G(u_k)) = 0$, so the partition by the points $u_K < s_{K-1} < u_{K-1} < \cdots < s_1 < u_1$
gives a sum at least $\sum_{k=1}^{K-1}s_k$, which tends to $\infty$ with $K$.
\end{solution}

\begin{exercise}\label{lebesgue-differentiation:exercise-strictly-singular}
Let $c$ be the Cantor function, extended by $c(t) = 0$ for $t < 0$ and $c(t) = 1$ for $t > 1$, and let $q_1, q_2, \ldots$ be an enumeration of the rationals in $(0, 1)$. Show that
$F(x) = \sum_n2^{-n}c(x - q_n)$ is continuous and strictly increasing on $[0, 1]$, but $F^{\prime} = 0$ almost everywhere.
\end{exercise}

\begin{solution}
Each term is continuous and nondecreasing, and the series converges uniformly, so $F$ is continuous and nondecreasing. If $0 \leq x < y \leq 1$, choose $q_n \in (x, y)$; then $c(x - q_n) = 0$ while
$c(y - q_n) > 0$, because $c(t) \geq c(3^{-k}) = 2^{-k} > 0$ for $t \geq 3^{-k}$. So $F(y) - F(x) \geq 2^{-n}(c(y - q_n) - c(x - q_n)) > 0$. Each $c(x - q_n)$ has derivative $0$ almost everywhere, since $c^{\prime} = 0$
outside the Cantor set and $c$ is constant outside $[0, 1]$. By Theorem~\ref{lebesgue-differentiation:theorem-fubini-series}, $F^{\prime} = \sum_n2^{-n}\cdot0 = 0$ almost everywhere.
\end{solution}
